\chapter{Why rotation needs an inertia tensor}
\section{Torque comes from force and lever arm}
For a force $\vect F$ applied at displacement $\vect r_i$ from the centre of mass,
\begin{equation}
\vect\tau_i=\vect r_i\times\vect F_i.
\end{equation}
For $\vect a=(a_x,a_y,a_z)^T$ and $\vect b=(b_x,b_y,b_z)^T$,
\begin{equation}
\vect a\times\vect b=
\begin{bmatrix}a_yb_z-a_zb_y\\a_zb_x-a_xb_z\\a_xb_y-a_yb_x\end{bmatrix}.
\end{equation}
This operation is the vector cross product. Its magnitude is
$\|\vect a\|\|\vect b\|\sin\gamma$ and its direction follows the right-hand rule.
It is not component-by-component multiplication.

For a rotor at $(x_i,y_i,z_i)$ with upward thrust $(0,0,T_i)$,
\begin{equation}
\vect\tau_{T,i}=
\begin{bmatrix}y_iT_i\\-x_iT_i\\0\end{bmatrix}.
\end{equation}
A rotor far from the centre creates a greater roll/pitch torque for the same
thrust. Its vertical height $z_i$ does not appear for this parallel upward force.

\section{Build angular momentum from individual mass elements}
For a mass element $\dd m$ at body position $\vect\rho$, rigid-body rotation
produces velocity $\vect\omega\times\vect\rho$ relative to the centre. Its angular
momentum about the centre is
\begin{align}
\dd\vect H&=\vect\rho\times(\vect\omega\times\vect\rho)\dd m\\
&=\bigl[(\vect\rho^T\vect\rho)I-\vect\rho\vect\rho^T\bigr]\vect\omega\dd m.
\end{align}
Integrate over the entire assembled body:
\begin{align}
\vect H_B&=J\vect\omega,\\
J&=\int\left(\|\vect\rho\|^2I-\vect\rho\vect\rho^T\right)\dd m\\
&=\begin{bmatrix}
\int(y^2+z^2)\dd m&-\int xy\dd m&-\int xz\dd m\\
-\int xy\dd m&\int(x^2+z^2)\dd m&-\int yz\dd m\\
-\int xz\dd m&-\int yz\dd m&\int(x^2+y^2)\dd m
\end{bmatrix}.
\end{align}
$J_x$ therefore measures resistance to angular acceleration around body x.
It depends on where mass is located, not only the total mass. Moving a battery
away from an axis increases inertia about that axis.

\section{Derive the rotating-frame derivative}
Let a physical vector have body components $\vect H_B$ and world components
$\vect H_W=R\vect H_B$. Differentiate with the product rule:
\begin{align}
\dot{\vect H}_W&=\dot R\vect H_B+R\dot{\vect H}_B.
\end{align}
Rotation kinematics gives $\dot R=R\crossmat{\vect\omega}$, where
\begin{equation}
\crossmat{\vect\omega}=
\begin{bmatrix}0&-r_b&q\\r_b&0&-p\\-q&p&0\end{bmatrix},
\qquad \crossmat{\vect\omega}\vect H_B=\vect\omega\times\vect H_B.
\end{equation}
Thus
\begin{align}
\dot{\vect H}_W
&=R\left(\dot{\vect H}_B+\vect\omega\times\vect H_B\right).
\end{align}
Angular-momentum balance says $\dot{\vect H}_W=\vect\tau_W$.
Multiply by $R^T$:
\begin{align}
\vect\tau_B&=\dot{\vect H}_B+\vect\omega\times\vect H_B.
\end{align}
For constant inertia, $\dot{\vect H}_B=J\dot{\vect\omega}$. Substitute
$\vect H_B=J\vect\omega$:
\begin{align}
\vect\tau_B&=J\dot{\vect\omega}+\vect\omega\times(J\vect\omega),\\
J\dot{\vect\omega}
&=\vect\tau_B-\vect\omega\times(J\vect\omega).
\end{align}
The cross term accounts for the body axes themselves rotating while the angular
momentum vector is described in those axes. It is not an extra arbitrary torque.

\section{Expand the nonlinear coupling}
If the selected body axes are principal axes,
$J=\diag(J_x,J_y,J_z)$ and
\begin{align}
J\vect\omega&=(J_xp,J_yq,J_zr_b)^T,\\
\vect\omega\times(J\vect\omega)
&=\begin{bmatrix}
(J_z-J_y)qr_b\\(J_x-J_z)pr_b\\(J_y-J_x)pq
\end{bmatrix}.
\end{align}
Subtract each component and divide by its inertia:
\begin{align}
\dot p&=\frac{\tau_x+(J_y-J_z)qr_b}{J_x},\\
\dot q&=\frac{\tau_y+(J_z-J_x)r_bp}{J_y},\\
\dot r_b&=\frac{\tau_z+(J_x-J_y)pq}{J_z}.
\end{align}
For a non-diagonal measured tensor, use the full matrix inverse or solve
$J\dot{\vect\omega}=\cdots$; do not delete products of inertia just for convenience.

\section{Angle derivatives are not generally gyro components}
For the stated Euler convention,
\begin{equation}
\begin{bmatrix}\dot\phi\\\dot\theta\\\dot\psi\end{bmatrix}=
\begin{bmatrix}
1&\sin\phi\tan\theta&\cos\phi\tan\theta\\
0&\cos\phi&-\sin\phi\\
0&\sin\phi/\cos\theta&\cos\phi/\cos\theta
\end{bmatrix}
\begin{bmatrix}p\\q\\r_b\end{bmatrix}.
\end{equation}
Only near modest attitudes/rates can $\dot\phi\approx p$ and $\dot\theta\approx q$
be substituted conveniently. The current independent angle/bias filters use that
approximation; the physical rotational model need not.

To avoid the Euler singularity at $\cos\theta=0$, propagate a unit quaternion.
For body-to-world quaternion $q=(q_w,q_x,q_y,q_z)$,
\begin{equation}
\dot q=\frac12 q\otimes(0,p,q_{\rm rate},r_b),\qquad \|q\|=1.
\end{equation}
Here $q_{\rm rate}$ means the gyro y rate, not the quaternion. A numerical
implementation must keep these two meanings distinct and normalize as needed.

\chapter{The four rotors and the nonlinear actuator model}
\section{Geometry and force allocation}
For motor diagonal $D$ in an X layout, $a=D/(2\sqrt2)$:
\begin{equation}
(x_i,y_i)=
(a,a),(a,-a),(-a,-a),(-a,a)
\end{equation}
in order FL, FR, RR, RL.
\begin{figure}[ht]
\centering
\begin{tikzpicture}[>=Latex,scale=1.1]
\draw[very thick,red!70!black] (0,0)--(-1.4,1.4) (0,0)--(1.4,1.4);
\draw[very thick,gray] (0,0)--(-1.4,-1.4) (0,0)--(1.4,-1.4);
\foreach \x/\y/\name in {-1.4/1.4/FL,1.4/1.4/FR,1.4/-1.4/RR,-1.4/-1.4/RL}{
\draw[fill=soft] (\x,\y) circle (.32); \node at (\x,\y){\name};}
\draw[->,thick] (0,0)--(0,2.1) node[above]{$x_B$ forward};
\draw[->,thick] (0,0)--(-2.1,0) node[left]{$y_B$ left};
\fill (0,0) circle (.06);
\node[right] at (.2,0){COM};
\node at (0,-2.1){Top view: $z_B$ points out of the page};
\end{tikzpicture}
\caption{Logical geometry; verify actual wiring, arm orientation and propeller
spin before assigning a physical motor number.}
\end{figure}
With reaction torque $Q_i=\kappa T_i$ and the selected yaw signs:
\begin{equation}
\begin{bmatrix}T\\\tau_x\\\tau_y\\\tau_z\end{bmatrix}=
\begin{bmatrix}
1&1&1&1\\a&-a&-a&a\\-a&-a&a&a\\-\kappa&\kappa&-\kappa&\kappa
\end{bmatrix}
\begin{bmatrix}T_{FL}\\T_{FR}\\T_{RR}\\T_{RL}\end{bmatrix}.
\end{equation}
This allocation is a geometric consequence of $\vect r\times\vect F$ plus yaw
reaction signs. A physically measured $Q_i(\Omega_i)$ can replace the constant
$\kappa$ model.

\section{Retain speed-squared thrust}
A common aerodynamic form is
\begin{equation}
T_i=k_T\Omega_i^2,\qquad Q_i=k_Q\Omega_i^2.
\end{equation}
More generally $k_T,k_Q$ depend on inflow, propeller condition, density and advance
ratio. These functions are not identified from a 9047 label alone.

An effective motor/ESC model can be written
\begin{align}
L\dot I_i&=V_{{eff},i}-R_m I_i-k_e\Omega_i,\\
J_r\dot\Omega_i&=k_\tau I_i-Q_i(\Omega_i)-B_r\Omega_i,\\
V_{{eff},i}&=\mu(u_i,V_{\rm pack},\hbox{ESC settings})V_{\rm pack}.
\end{align}
The ESC is an electronic commutator, so these are effective averaged equations;
phase/line voltage conventions and efficiency must be specified before interpreting
$k_e,k_\tau$. Do not silently equate a PWM input percentage with an effective
phase voltage or a thrust percentage.

The KV rating indicates an approximate no-load speed/voltage relationship.
It does not supply $R_m,L,J_r,k_T,k_Q$, loaded RPM, ESC map $\mu$, or transient
delay. We therefore retain these coefficients symbolically until measurement.

\section{Rotor angular momentum: optional additional physics}
If $\sigma_i=\pm1$ denotes rotor spin direction,
\begin{equation}
\vect H_{\rm rot}=\sum_i\sigma_iJ_r\Omega_i\vect e_3.
\end{equation}
For an assembled-airframe inertia $J$ defined for frozen relative rotor angles,
the total rotating-body momentum balance can be expanded as
\begin{equation}
J\dot{\vect\omega}
=\vect\tau_{\rm external}
-\vect\omega\times(J\vect\omega)
-\dot{\vect H}_{\rm rot}
-\vect\omega\times\vect H_{\rm rot}.
\end{equation}
External rotor aerodynamic drag has sign opposite $\sigma_i$; if using the
usual body reaction-torque sign $s_i$, then $s_i=-\sigma_i$.
Do not count the same motor/rotor torque twice by adding it both as an external
reaction and as an internal rotor acceleration contribution.

These extra terms need rotor inertia/spin information. They are model extensions,
not claims that the present Gazebo adapter has identified or enabled them.

\section{Battery and actuator coupling}
A symbolic electrical extension is
\begin{align}
V_{\rm term}&=V_{\rm oc}(z,T_b)-I_{\rm pack}R_0-V_p,\\
\dot V_p&=-\frac{V_p}{R_1C_1}+\frac{I_{\rm pack}}{C_1},\\
\dot z&=-\frac{\eta(I_{\rm pack},T_b)I_{\rm pack}}{3600C_{\rm eff}},
\qquad z\in[0,1].
\end{align}
$C_{\rm eff}$ is in Ah and positive current means discharge. Phase motor currents
are not automatically identical to pack current; inverter efficiency and duty
matter. Voltage sag changes available speed/thrust, connecting electrical and
mechanical dynamics.

\important{For current reproducible flight results, the adapter still uses
declared first-order thrust lag and normalized linear thrust, while Gazebo
integrates nonlinear rigid-body/contact dynamics. The richer equations above
define the model to identify; no arbitrary electrical coefficients were invented
or silently substituted into an already-recorded flight.}

\section{Why a single transfer function cannot represent all of this}
Transfer functions describe linear time-invariant input/output models with
specified initial conditions. Our exact equations contain $\Omega^2$,
$T\cos\phi\cos\theta$, and $pq$ products. In general,
\begin{equation}
\mathcal L\{\Omega(t)^2\}\ne\bigl(\mathcal L\{\Omega(t)\}\bigr)^2.
\end{equation}
We use the nonlinear equations for prediction and simulation, and a local
Jacobian model for transfer-function/gain design. Multiple operating points or
nonlinear control methods can extend this design later; pretending the full
nonlinear aircraft has one global scalar transfer function would be incorrect.
